% Figura redibujada a partir de "Imagenes/Tablas de MOSFET.PNG" (Tema 5).
% Coordenadas en píxeles del PNG original (1289 x 775 px), escaladas para que
% ocupe lo mismo que el PNG en el documento: 483.7 pt x 290.8 pt (\linewidth).
\documentclass[12pt,border=0pt]{standalone}
\usepackage[spanish,es-tabla]{babel}
\usepackage[utf8]{inputenc}
\usepackage{amsmath}
\usepackage{xcolor}
\usepackage{tikz}
\usepackage{pgfplots}
\usepackage{circuitikz}
\usetikzlibrary{babel, circuits.ee}
\usetikzlibrary{decorations.pathreplacing}

\definecolor{applegreen}{rgb}{0.55, 0.71, 0.0}
\definecolor{airforceblue}{rgb}{0.36, 0.54, 0.66}
\definecolor{amber}{rgb}{1.0, 0.49, 0.0}

\newcommand{\diff}[3][1]{%
	\IfBlankTF{#1}{\frac{\mathrm{d}{#2}}{\mathrm{d}{#3}}}%
	{\frac{\mathrm{d}^{#1} {#2}}{\mathrm{d}{#3}^{#1}}}%
}%
\newcommand{\pdiff}[3][1]{%
	\IfBlankTF{#1}{\frac{\partial {#2}}{\partial {#3}}}%
	{\frac{\partial^{#1} {#2}}{\partial {#3}^{#1}}}%
}%

% Símbolo de MOSFET de 4 terminales (con sustrato B).
% #1: posición de la puerta; #2: 1 = empobrecimiento (canal continuo);
% #3: 1 = canal P (flecha del sustrato hacia fuera).
\newcommand{\mosfetB}[3]{%
    \begin{scope}[shift={#1}]
        \draw (-31,30) -- (0,30) (0,-28) -- (0,30);
        \ifnum#2=1 \draw (15,-48) -- (15,46);
        \else \draw (15,-48) -- (15,-30) (15,-10) -- (15,10) (15,28) -- (15,46);
        \fi
        \draw (15,-40) -- (30,-40) -- (30,-55);
        \draw (15,38) -- (30,38) -- (30,52);
        \ifnum#3=1 \draw[-latex] (15,0) -- (50,0);
        \else \draw[-latex] (50,0) -- (16,0);
        \fi
        \node at (-40,0) {G};
        \node at (67,-45) {D};
        \node at (65,-2) {B};
        \node at (67,38) {S};
    \end{scope}}
% Símbolo de MOSFET de 3 terminales.
% #1: posición de la puerta; #2: 1 = empobrecimiento (canal grueso);
% #3: 1 = canal P (flecha de la fuente hacia el canal).
\newcommand{\mosfet}[3]{%
    \begin{scope}[shift={#1}]
        \draw (-30,30) -- (0,30) (0,-30) -- (0,30);
        \draw (15,-50) -- (15,40);
        \ifnum#2=1 \fill (16,-38) rectangle (23,30); \fi
        \draw (15,-43) -- (47,-43) -- (47,-56);
        \ifnum#3=1 \draw[-latex] (47,23) -- (16,23);
        \else \draw[-latex] (15,23) -- (47,23);
        \fi
        \draw (47,23) -- (47,40);
        \node at (-41,3) {G};
        \node at (65,-45) {D};
        \node at (65,35) {S};
    \end{scope}}
% Curvas de salida I_D-V_DS. #1: origen; #2: sentido (1 canal N, -1 canal P);
% #3: lista de nivel/codo/final (en píxeles, respecto al origen).
\newcommand{\curvassalida}[3]{%
    \begin{scope}[shift={#1}, xscale=#2, yscale=#2]
        \foreach \nivel/\codo/\fin in {#3}
            \draw[curva] (0,0) to[out=72, in=180] (\codo,-\nivel) -- (\fin,-\nivel);
    \end{scope}}

\begin{document}
\begin{tikzpicture}[x=0.37525pt, y=-0.37525pt, font=\scriptsize,
    eje/.style={-latex}, curva/.style={blue, thick}]
    \useasboundingbox (0,0) rectangle (1289,775);
    % ---- Rejilla de la tabla ----
    \foreach \x in {203,383,712,940,1150} \draw (\x,8) -- (\x,770);
    \foreach \y in {8,97,268,433,603,770} \draw (203,\y) -- (1150,\y);
    % ---- Cabeceras ----
    \begin{scope}[font=\footnotesize, align=center]
        \node at (293,48) {Tipo};
        \node at (547,48) {Símbolos};
        \node at (826,48) {Características\\de salida};
        \node at (1045,48) {Características de\\transferencia};
    \end{scope}
    \begin{scope}[font=\scriptsize\bfseries, align=center]
        \node at (98,56) {Condición\\para conducción};
        \node at (1220,56) {Potencial\\umbral};
    \end{scope}
    % ---- Columna izquierda: condición de conducción ----
    \draw[decorate, decoration={brace, amplitude=7pt}] (180,430) -- (180,105);
    \draw[decorate, decoration={brace, amplitude=7pt}] (180,765) -- (180,440);
    \node[font=\footnotesize] at (80,272) {$V_\mathrm{GS} > V_\mathrm{T}$};
    \node[font=\footnotesize] at (78,600) {$V_\mathrm{GS} < V_\mathrm{T}$};
    % ---- Columna derecha: tensión umbral ----
    \begin{scope}[font=\footnotesize]
        \node at (1228,178) {$V_\mathrm{T} > 0$};
        \node at (1228,345) {$V_\mathrm{T} < 0$};
        \node at (1228,527) {$V_\mathrm{T} < 0$};
        \node at (1228,690) {$V_\mathrm{T} > 0$};
    \end{scope}
    % ---- Tipo ----
    \begin{scope}[align=center]
        \node at (293,175) {Canal N\\Enriquecimiento\\Norm. OFF\\Realce};
        \node at (293,337) {Canal N\\Empobrecimiento\\Norm. ON\\Vaciamiento};
        \node at (293,506) {Canal P\\Enriquecimiento\\Norm. OFF\\Realce};
        \node at (293,668) {Canal P\\Empobrecimiento\\Norm. ON\\Vaciamiento};
    \end{scope}
    % ---- Símbolos ----
    \begin{scope}[]
        \mosfetB{(442,172)}{0}{0}  \mosfet{(604,172)}{0}{0}
        \mosfetB{(442,340)}{1}{0}  \mosfet{(604,340)}{1}{0}
        \mosfetB{(442,520)}{0}{1}  \mosfet{(604,520)}{0}{1}
        \mosfetB{(442,680)}{1}{1}  \mosfet{(592,680)}{1}{1}
    \end{scope}
    % ---- Características de salida ----
    % Canal N, enriquecimiento
    \draw[eje] (742,222) -- (742,105);
    \draw[eje] (735,215) -- (890,215);
    \curvassalida{(742,215)}{1}{93/70/113, 57/55/98, 33/42/98, 15/32/98}
    \node at (764,120) {$I_\mathrm{D}$};
    \node at (818,232) {$V_\mathrm{DS}$};
    \node at (862,156) {$V_\mathrm{GS}$};
    \draw[-latex] (890,145) -- (890,117);
    \node at (893,107) {$+$};
    \draw[-latex] (876,236) -- (900,236);
    \node at (915,232) {$+$};
    % Canal N, empobrecimiento
    \draw[eje] (742,397) -- (742,280);
    \draw[eje] (735,390) -- (890,390);
    \curvassalida{(742,390)}{1}{97/72/113, 55/55/98, 30/42/98, 12/32/98}
    \node at (764,295) {$I_\mathrm{D}$};
    \node at (818,408) {$V_\mathrm{DS}$};
    \node at (872,332) {$V_\mathrm{GS}\!=\!0$};
    \draw[-latex] (890,320) -- (890,292);
    \node at (893,282) {$+$};
    \draw[-latex] (876,410) -- (900,410);
    \node at (915,407) {$+$};
    % Canal P, enriquecimiento
    \draw[eje] (750,468) -- (912,468);
    \draw (893,462) -- (893,592);
    \curvassalida{(893,468)}{-1}{27/45/115, 47/60/120, 67/75/122, 90/90/118}
    \node at (802,454) {$V_\mathrm{DS}$};
    \draw[-latex] (785,452) -- (757,452);
    \node at (745,463) {$-$};
    \node at (770,515) {$V_\mathrm{GS}$};
    \draw[-latex] (762,528) -- (762,560);
    \node at (762,573) {$-$};
    \node at (870,556) {$I_\mathrm{D}$};
    % Canal P, empobrecimiento
    \draw[eje] (760,622) -- (918,622);
    \draw (895,616) -- (895,745);
    \curvassalida{(895,622)}{-1}{28/45/120, 48/60/125, 70/75/128, 100/92/130}
    \node at (806,610) {$V_\mathrm{DS}$};
    \node at (733,635) {$-$};
    \node at (765,672) {$V_\mathrm{GS}\!=\!0$};
    \draw[-latex] (752,688) -- (752,722);
    \node at (752,736) {$-$};
    \node at (868,718) {$I_\mathrm{D}$};
    % ---- Características de transferencia ----
    % Canal N, enriquecimiento
    \draw[eje] (945,218) -- (1122,218);
    \draw[eje] (1027,218) -- (1027,116);
    \draw[curva] (1060,218) .. controls (1080,216) and (1095,160) .. (1118,130);
    \node at (1045,128) {$I_\mathrm{D}$};
    \node at (1106,190) {$V_\mathrm{GS}$};
    \node at (1042,234) {$V_\mathrm{T}$};
    \node at (956,232) {$-$};
    \node at (1113,234) {$+$};
    % Canal N, empobrecimiento
    \draw[eje] (948,385) -- (1127,385);
    \draw[eje] (1027,385) -- (1027,283);
    \draw[curva] (985,385) .. controls (1010,383) and (1030,340) .. (1088,315);
    \node at (1052,293) {$I_\mathrm{D}$};
    \node at (1112,358) {$V_\mathrm{GS}$};
    \node at (956,400) {$-$};
    \node at (1118,400) {$+$};
    % Canal P, enriquecimiento
    \draw[eje] (962,485) -- (1142,485);
    \draw[eje] (1042,592) -- (1042,455);
    \draw[curva] (965,570) .. controls (985,560) and (1000,500) .. (1017,485);
    \node at (1002,465) {$V_\mathrm{T}$};
    \node at (1068,469) {$I_\mathrm{D}$};
    \node at (1128,467) {$+$};
    \node at (963,480) {$-$};
    \node at (1118,513) {$V_\mathrm{GS}$};
    % Canal P, empobrecimiento
    \draw[eje] (950,650) -- (1132,650);
    \draw[eje] (1030,735) -- (1030,623);
    \draw[curva] (965,725) .. controls (1000,705) and (1020,655) .. (1060,652) -- (1080,651);
    \node at (1062,632) {$I_\mathrm{D}$};
    \node at (1122,634) {$+$};
    \node at (963,645) {$-$};
    \node at (1105,678) {$V_\mathrm{GS}$};
\end{tikzpicture}
\end{document}
